\documentclass{stud-conf-pmk}

\begin{document}
\year{2026}

\begin{paper}
\author[Гаврилов А.\,Д.]{Гаврилов Антон Дмитриевич}
\title{Детекция и трекинг игроков футбольного матча с~использованием методов компьютерного зрения}
\speciality{Прикладная математика и информатика}
\student{Студент магистратуры}
\presented{Представил доцент кафедры информатики М.\,Ю.\,Кудряшов}
\UDC{004.932.2, 004.85}
\abstract{В работе рассматривается задача автоматической детекции и~трекинга игроков футбольного матча на~видеозаписях. Предложена модульная система, включающая детектор объектов YOLOv8m, дообученный на~специализированном датасете, и~алгоритм многообъектного трекинга ByteTrack. Проведена экспериментальная оценка на~видеофрагменте реального матча. Достигнуты значения mAP50~=~0{,}831 и~Precision~=~0{,}960. Исследована возможность классификации игроков по~командам методом K-Means.}
\keywords{детекция объектов, YOLOv8, многообъектный трекинг, ByteTrack, компьютерное зрение, футбол}
\maketitle

\introduction

Профессиональный футбол генерирует значительные объёмы видеоданных: каждый матч фиксируется множеством камер. Тренерские штабы, аналитические компании и~вещательные организации нуждаются в~инструментах автоматического анализа происходящего на~поле. Ручная обработка видеозаписей требует значительных временных затрат и~подвержена субъективности~\cite{bib_sports_cv}.

Задача автоматического анализа футбольного видео включает несколько подзадач: детекцию (обнаружение) объектов~--- игроков, вратарей, судей и~мяча; трекинг~--- отслеживание каждого объекта на~протяжении видеопоследовательности с~сохранением уникального идентификатора; и~классификацию~--- определение принадлежности игрока к~команде.

Развитие глубокого обучения, в~частности свёрточных нейронных сетей~\cite{bib_goodfellow}, привело к~появлению эффективных одностадийных детекторов семейства YOLO~\cite{bib_yolo_original}, способных работать в~реальном времени. В~сочетании с~современными алгоритмами многообъектного трекинга, такими как ByteTrack~\cite{bib_bytetrack}, это открывает возможности для~построения практических систем анализа спортивного видео.

Целью данной работы является разработка системы автоматической детекции и~трекинга игроков футбольного матча на~видеозаписях на~основе методов глубокого обучения и~компьютерного зрения.


\section{Методы детекции и~трекинга}

Для~детекции объектов выбрана архитектура YOLOv8~\cite{bib_yolov8_ultralytics}~--- восьмая версия семейства YOLO. В~отличие от~двухстадийных детекторов, таких как Faster~R-CNN~\cite{bib_faster_rcnn}, YOLOv8 является одностадийным: за~один проход через нейронную сеть формируются все предсказания координат ограничивающих рамок и~классов объектов. Это обеспечивает высокую скорость работы при~сохранении точности.

Архитектура YOLOv8 состоит из~трёх основных компонентов: Backbone (CSPDarknet) для~извлечения визуальных признаков из~изображения; Neck (FPN~+~PAN) для~объединения признаков на~разных масштабах; Head (Decoupled Head)~--- разделённая голова, в~которой задачи классификации и~регрессии координат решаются независимо. Ключевой особенностью YOLOv8 является отсутствие предопределённых якорных рамок (anchor-free подход), что~упрощает настройку модели.

Используется модель YOLOv8m (medium), обеспечивающая баланс между~точностью и~вычислительной эффективностью. Детекция ведётся по~четырём классам: \texttt{player}, \texttt{goalkeeper}, \texttt{referee}, \texttt{ball}.

Для~многообъектного трекинга применяется алгоритм ByteTrack~\cite{bib_bytetrack}. В~его основе лежит фильтр Калмана~\cite{bib_kalman}, предсказывающий положение объекта на~следующем кадре по~текущему состоянию. Ключевая особенность ByteTrack~--- двухуровневое сопоставление детекций с~треками. На~первом этапе уверенные детекции (с~высоким значением confidence) сопоставляются с~предсказанными положениями треков на~основе метрики IoU (Intersection over Union). На~втором этапе повторно обрабатываются оставшиеся, менее уверенные детекции. Такой подход позволяет сохранять идентичность объектов при~окклюзиях~--- ситуациях, когда один объект частично или полностью перекрывает другой.


\section{Датасет и~обучение модели}

Для~дообучения модели используется датасет football-players-detection с~платформы Roboflow~\cite{bib_roboflow_dataset}, содержащий 255~размеченных изображений футбольных матчей. Разметка включает четыре класса объектов. Датасет разделён на~обучающую (204~изображения), валидационную (38) и~тестовую (13) выборки.

Обучение проводится методом transfer learning (перенос обучения): модель YOLOv8m, предобученная на~датасете COCO (330\,000 изображений, 80~классов), дообучается на~футбольном датасете. Данный подход позволяет эффективно использовать выученные низкоуровневые визуальные признаки и~достигать высокого качества при~ограниченном объёме обучающих данных.

Обучение проведено в~среде Google Colab на~GPU NVIDIA Tesla T4 (16\,ГБ видеопамяти) в~течение 200~эпох при~входном разрешении $640 \times 640$~пикселей. Для~повышения обобщающей способности модели применяются аугментации: mosaic (объединение четырёх изображений), mixup (наложение двух изображений), случайные аффинные преобразования и~модификации в~цветовом пространстве HSV.


\section{Результаты экспериментов}

Качество детекции оценивается стандартными метриками. Метрика IoU (Intersection over Union) определяется как отношение площади пересечения предсказанной и~эталонной рамок к~площади их объединения:
$$\text{IoU} = \frac{|A \cap B|}{|A \cup B|},$$
где $A$~--- предсказанная ограничивающая рамка, $B$~--- эталонная разметка.

На~основе IoU вычисляются Precision (доля верных детекций среди~всех предсказанных), Recall (доля найденных объектов среди~всех реальных) и~mAP (mean Average Precision)~--- площадь под~кривой Precision-Recall, усреднённая по~всем классам.

Результаты на~валидационной выборке представлены в~табл.~\ref{tab:metrics}.

\begin{table}
\centering
\caption{Метрики качества детекции YOLOv8m.}
\begin{tabular}{|l|l|l|l|l|}
\hline
Класс & Precision & Recall & mAP50 & mAP50-95 \\
\hline
player     & 0{,}974 & 0{,}981 & 0{,}982 & 0{,}766 \\
\hline
goalkeeper & 0{,}965 & 0{,}776 & 0{,}969 & 0{,}713 \\
\hline
referee    & 0{,}952 & 0{,}942 & 0{,}925 & 0{,}644 \\
\hline
ball       & 0{,}948 & 0{,}400 & 0{,}448 & 0{,}173 \\
\hline
Среднее    & 0{,}960 & 0{,}775 & 0{,}831 & 0{,}574 \\
\hline
\end{tabular}
\label{tab:metrics}
\end{table}

Наилучшие результаты достигнуты для~класса \texttt{player}: mAP50~=~0{,}982 на~754~экземплярах. Класс \texttt{ball} показывает наихудший результат (mAP50-95~=~0{,}173), что~объясняется малым размером объекта (порядка $10 \times 10$~пикселей при~входном разрешении) и~размытием при~быстром движении. Скорость инференса на~Tesla~T4 составляет примерно 19~FPS (8{,}4~мс препроцессинг, 39{,}5~мс детекция, 4{,}7~мс постпроцессинг).

Тестирование полного пайплайна (детекция~+ трекинг~+ визуализация) проведено на~30-секундном видеофрагменте матча Бельгия\,--\,Россия (EURO~2020), содержащем 750~кадров при~25~FPS. Система стабильно обнаруживает 23~объекта на~кадре: 22~полевых игрока и~мяч. Алгоритм ByteTrack обеспечивает устойчивое сохранение уникальных идентификаторов объектов между~кадрами. Потеря идентичности наблюдается при~длительных окклюзиях (более~1~секунды).

Дополнительно исследована возможность классификации игроков по~командам на~основе цвета формы. Из~ограничивающей рамки каждого игрока выделяется область торса, выполняется фильтрация зелёных пикселей (маска в~пространстве HSV), вычисляется медианный цвет оставшихся пикселей, после чего~применяется алгоритм K-Means~\cite{bib_kmeans} с~$k=2$ для~разделения на~две группы. Эксперименты показали, что~данный подход неустойчив при~низком разрешении видео: область торса содержит слишком мало пикселей, тени и~блики искажают цвет, кластеризация нестабильна. Модуль не~включён в~финальную версию системы и~рассматривается как~направление дальнейшего развития.

\conclusion

В~работе разработана модульная система автоматической детекции и~трекинга игроков футбольного матча на~видеозаписях. Модель YOLOv8m дообучена на~специализированном датасете и~достигла значения mAP50~=~0{,}831 при~Precision~=~0{,}960. Реализован полный пайплайн обработки видео с~использованием алгоритма ByteTrack для~многообъектного трекинга. Экспериментальная оценка на~реальном видео матча подтвердила работоспособность системы.

В~качестве направлений дальнейшего развития можно выделить: использование re-identification сетей для~классификации команд; обработку видео более~высокого разрешения для~улучшения детекции мяча; построение тепловых карт перемещений игроков на~основе данных трекинга.

\begin{thebibliography}{99}

\bibitem{bib_goodfellow}
Goodfellow, I. Deep Learning~/ I.\,Goodfellow, Y.\,Bengio, A.\,Courville.~--- Cambridge~: The MIT Press, 2016.~--- 800\,p.~--- URL: \url{https://www.deeplearningbook.org/}

\bibitem{bib_resnet}
He, K. Deep Residual Learning for Image Recognition~/ K.\,He, X.\,Zhang, S.\,Ren, J.\,Sun~// Proceedings of the IEEE Conference on Computer Vision and Pattern Recognition (CVPR).~--- 2016.~--- P.\,770--778.~--- DOI: \url{https://doi.org/10.1109/CVPR.2016.90}

\bibitem{bib_yolov8_ultralytics}
Jocher, G. Ultralytics YOLOv8~/ G.\,Jocher, A.\,Chaurasia, J.\,Qiu.~--- 2023.~--- URL: \url{https://github.com/ultralytics/ultralytics} (дата обращения: 01.12.2025).

\bibitem{bib_kalman}
Kalman, R.\,E. A New Approach to Linear Filtering and Prediction Problems~// Journal of Basic Engineering.~--- 1960.~--- Vol.\,82, №\,1.~--- P.\,35--45.~--- DOI: \url{https://doi.org/10.1115/1.3662552}

\bibitem{bib_kmeans}
Lloyd, S.\,P. Least Squares Quantization in PCM~// IEEE Transactions on Information Theory.~--- 1982.~--- Vol.\,28, №\,2.~--- P.\,129--137.~--- DOI: \url{https://doi.org/10.1109/TIT.1982.1056489}

\bibitem{bib_yolo_original}
Redmon, J. You Only Look Once: Unified, Real-Time Object Detection~/ J.\,Redmon, S.\,Divvala, R.\,Girshick, A.\,Farhadi~// Proceedings of the IEEE Conference on Computer Vision and Pattern Recognition (CVPR).~--- 2016.~--- P.\,779--788.~--- DOI: \url{https://doi.org/10.1109/CVPR.2016.91}

\bibitem{bib_faster_rcnn}
Ren, S. Faster R-CNN: Towards Real-Time Object Detection with Region Proposal Networks~/ S.\,Ren, K.\,He, R.\,Girshick, J.\,Sun~// Advances in Neural Information Processing Systems (NeurIPS).~--- 2015.~--- Vol.\,28.

\bibitem{bib_roboflow_dataset}
Football Players Detection Dataset~// Roboflow Universe~: [сайт].~--- URL: \url{https://universe.roboflow.com/roboflow-jvuqo/football-players-detection-3zvbc} (дата обращения: 01.12.2025).~--- Загл. с~титул. экрана.

\bibitem{bib_sports_cv}
Thomas, G. Computer Vision for Sports: Current Applications and Research Topics~/ G.\,Thomas, R.\,Gade, T.\,B.\,Moeslund [et~al.]~// Computer Vision and Image Understanding.~--- 2017.~--- Vol.\,159.~--- P.\,3--18.~--- DOI: \url{https://doi.org/10.1016/j.cviu.2017.04.011}

\bibitem{bib_bytetrack}
Zhang, Y. ByteTrack: Multi-Object Tracking by Associating Every Detection Box~/ Y.\,Zhang, P.\,Sun, Y.\,Jiang [et~al.]~// Proceedings of the European Conference on Computer Vision (ECCV).~--- 2022.~--- P.\,1--21.~--- DOI: \url{https://doi.org/10.1007/978-3-031-20047-2_1}

\end{thebibliography}

\end{paper}
\end{document}
